\hypertarget{ip-address-manipulation-and-rpc-ipaddress-xmlrpc}{%
\chapter{\texorpdfstring{IP Address Manipulation and RPC
(\texttt{ipaddress},
\texttt{xmlrpc})}{IP Address Manipulation and RPC (ipaddress, xmlrpc)}}\label{ip-address-manipulation-and-rpc-ipaddress-xmlrpc}}

\hypertarget{ipaddress-vectorized-network-math}{%
\subsection{\texorpdfstring{56.1 \texttt{ipaddress}: Vectorized Network
Math}{56.1 ipaddress: Vectorized Network Math}}\label{ipaddress-vectorized-network-math}}

Manipulating IP ranges with regex is a recipe for security
vulnerabilities. \passthrough{\lstinline!ipaddress!} provides objects
for IPv4 and IPv6 addresses and networks.

\hypertarget{internal-representations}{%
\subsubsection{1. Internal
Representations}\label{internal-representations}}

\begin{itemize}
\tightlist
\item
  \textbf{IPv4}: Stored as a 32-bit Python
  \passthrough{\lstinline!int!}.
\item
  \textbf{IPv6}: Stored as a 128-bit Python
  \passthrough{\lstinline!int!}.
\item
  \textbf{Performance}: Operations like
  \passthrough{\lstinline!addr in network!} are implemented using fast
  bitwise mask operations
  (\passthrough{\lstinline!(addr\_int \& mask) == network\_int!}),
  making them extremely efficient for high-speed firewall log analysis.
\end{itemize}

\hypertarget{xmlrpc-simple-remote-procedure-calls}{%
\subsection{\texorpdfstring{56.2 \texttt{xmlrpc}: Simple Remote
Procedure
Calls}{56.2 xmlrpc: Simple Remote Procedure Calls}}\label{xmlrpc-simple-remote-procedure-calls}}

XML-RPC is a legacy but still widely used protocol for calling functions
across the network. * \textbf{\passthrough{\lstinline!ServerProxy!}}:
Uses Python's \passthrough{\lstinline!\_\_getattr\_\_!} dunder method to
dynamically map local method calls to remote network requests. *
\textbf{Serialization}: It uses the \passthrough{\lstinline!xml.etree!}
module to convert Python types (ints, dicts, lists) into the XML format
required by the protocol.

\begin{center}\rule{0.5\linewidth}{0.5pt}\end{center}
